%! Algebra 2 — Unit 1, Lesson 3 — Exit Ticket (Key)
\documentclass[10pt]{article}
\usepackage{algebra2-article}
\usepackage{algebra2-boxes}
\usepackage{algebra2-key}

%=============================================================================
\begin{document}

\pageheader{Unit 1, Lesson 3}{Exit Ticket --- Key}
\namedateperiod

\vspace{0.3in}

\begin{enumerate}[leftmargin=*, itemsep=1.8em]

    % ── Question 1 ──────────────────────────────────────────────────────────
    \item For $f(x) = x^2 - 2x - 3$, identify each feature.

    \vspace{0.1in}
    \ansline{Factor: $f(x) = (x-3)(x+1)$; zeros at $x = 3$ and $x = -1$}
    \ansline{Axis of symmetry: $x = -\dfrac{-2}{2(1)} = 1$}
    \ansline{Vertex: $f(1) = 1 - 2 - 3 = -4$; vertex $(1, -4)$}
    \ansline{$y$-intercept: $f(0) = -3$; point $(0, -3)$}

    \vspace{0.05in}
    Zeros: \ans{$x = 3,\; x = -1$} \quad Vertex: \ans{$(1, -4)$} \quad $y$-intercept: \ans{$(0,-3)$}

    \begin{teachernote}
        Common errors: (1) computing $-b/2a$ as $-(-2)/2 = 1$ but then forgetting to evaluate $f(1)$;
        (2) listing the $y$-intercept as $-3$ without the coordinate notation.
        Cold-call: ``How did you find the vertex?''
    \end{teachernote}

    % ── Question 2 ──────────────────────────────────────────────────────────
    \item $y = 4.1x + 22$ (months since January; new customers)

    \vspace{0.1in}
    \ans{Slope $4.1$: each month, the business gains approximately $4.1$ new customers (constant rate of increase).}\\[4pt]
    \ans{$y$-intercept $22$: in January (month $0$) the business had approximately $22$ new customers.}

    \begin{teachernote}
        Watch for students who say ``the slope is the starting value'' or confuse the two.
        Reinforce: slope = rate of change (per unit of $x$); $y$-intercept = value when $x = 0$.
    \end{teachernote}

    % ── Question 3 (SOL-style MC) ────────────────────────────────────────────
    \item Which description fits a \emph{linear} scatterplot?

    \vspace{0.2in}
    \begin{tabularx}{\linewidth}{X X}
        \textbf{A.}\quad Points curve upward, then level off
        & \ans{\textbf{B.}\quad Points are clustered near a straight diagonal line} \\[12pt]
        \textbf{C.}\quad Points form a symmetric U-shape
        & \textbf{D.}\quad Points are randomly scattered with no visible pattern
    \end{tabularx}

    \vspace{0.2in}
    \ans{\textbf{B} is correct.} A linear model fits data that clusters near a straight line.
    A is exponential/logarithmic; C is quadratic; D suggests no relationship.

    \begin{teachernote}
        \textbf{Distractor analysis:}
        \textbf{A} --- students may associate ``rising'' with linear; clarify that leveling-off is a curve, not a line.
        \textbf{C} --- students who confuse ``symmetric'' with ``straight.''
        \textbf{D} --- students who are unsure what ``no pattern'' looks like.
    \end{teachernote}

\end{enumerate}

\end{document}
